\documentclass[10pt,a4paper]{amsart}
\usepackage{amsmath,amssymb,amsthm,graphicx}
% Retained preamble for the Stacks Project chapter snapshot in this directory.
% The chapter's own preamble.tex is a copy of the shared Stacks Project preamble; the block
% below is sliced MECHANICALLY, line for line, from it. Nothing was added, removed or
% redefined; the wrapper class, the xr-hyper/hyperref wiring and the externaldocument list
% are omitted because the assembler supplies the document class.
%
% lines taken from preamble.tex: 3, 5, 9, 11, 12, 156, 157, 158, ... (51 lines total)
%
% package coverage restored from the chapter's own preamble (2026-10-13): xy, xy-2cell

\usepackage[all]{xy}
\xyoption{2cell}
\UseAllTwocells

\newenvironment{reference}{\comment}{\endcomment}
\newenvironment{slogan}{\comment}{\endcomment}
\newenvironment{history}{\comment}{\endcomment}
\theoremstyle{plain}
\newtheorem{theorem}[subsection]{Theorem}
\newtheorem{proposition}[subsection]{Proposition}
\newtheorem{lemma}[subsection]{Lemma}
\theoremstyle{definition}
\newtheorem{definition}[subsection]{Definition}
\newtheorem{example}[subsection]{Example}
\newtheorem{exercise}[subsection]{Exercise}
\newtheorem{situation}[subsection]{Situation}
\theoremstyle{remark}
\newtheorem{remark}[subsection]{Remark}
\newtheorem{remarks}[subsection]{Remarks}
\numberwithin{equation}{subsection}
\def\lim{\mathop{\mathrm{lim}}\nolimits}
\def\colim{\mathop{\mathrm{colim}}\nolimits}
\def\Spec{\mathop{\mathrm{Spec}}}
\def\Hom{\mathop{\mathrm{Hom}}\nolimits}
\def\Ext{\mathop{\mathrm{Ext}}\nolimits}
\def\SheafHom{\mathop{\mathcal{H}\!\mathit{om}}\nolimits}
\def\SheafExt{\mathop{\mathcal{E}\!\mathit{xt}}\nolimits}
\def\Sch{\mathit{Sch}}
\def\Mor{\mathop{\mathrm{Mor}}\nolimits}
\def\Ob{\mathop{\mathrm{Ob}}\nolimits}
\def\Sh{\mathop{\mathit{Sh}}\nolimits}
\def\NL{\mathop{N\!L}\nolimits}
\def\CH{\mathop{\mathrm{CH}}\nolimits}
\def\proetale{{pro\text{-}\acute{e}tale}}
\def\etale{{\acute{e}tale}}
\def\QCoh{\mathit{QCoh}}
\def\Ker{\mathop{\mathrm{Ker}}}
\def\Im{\mathop{\mathrm{Im}}}
\def\Coker{\mathop{\mathrm{Coker}}}
\def\Coim{\mathop{\mathrm{Coim}}}
\DeclareMathSymbol{\boxtimes}{\mathbin}{AMSa}{"02}
\def\QCohstack{\mathcal{QC}\!\mathit{oh}}
\def\Cohstack{\mathcal{C}\!\mathit{oh}}
\def\Spacesstack{\mathcal{S}\!\mathit{paces}}
\def\Quotfunctor{\mathrm{Quot}}
\def\Hilbfunctor{\mathrm{Hilb}}
\def\Curvesstack{\mathcal{C}\!\mathit{urves}}
\def\Polarizedstack{\mathcal{P}\!\mathit{olarized}}
\def\Complexesstack{\mathcal{C}\!\mathit{omplexes}}
\def\Pic{\mathop{\mathrm{Pic}}\nolimits}
\def\Picardstack{\mathcal{P}\!\mathit{ic}}
\def\Picardfunctor{\mathrm{Pic}}
\def\Deformationcategory{\mathcal{D}\!\mathit{ef}}
\usepackage{cleveref}
\crefname{originallemma}{original lemma}{original lemmas}
\Crefname{originallemma}{Original lemma}{Original lemmas}
\makeatletter
\newlabel{lemma-pullback-cotangent-morphism-topoi}{{08SV}{source}{}{sourcepointer-1}{}}
\newlabel{lemma-pullback-cotangent-morphism-topoi@cref}{{[originallemma][0][]08SV}{[1][0][]source}{}{}{}}
\makeatother
\crefname{originallemma}{original lemma}{original lemmas}
\Crefname{originallemma}{Original lemma}{Original lemmas}
\makeatletter
\newlabel{lemma-triangle-sheaves-rings}{{08SZ}{source}{}{sourcepointer-2}{}}
\newlabel{lemma-triangle-sheaves-rings@cref}{{[originallemma][0][]08SZ}{[1][0][]source}{}{}{}}
\makeatother
\title{The distinguished triangle for a composition of morphisms of ringed spaces}
\author{The Stacks Project Authors}
\date{stacks-project:cotangent.tex:local-snapshot-20261012; GNU FDL 1.2 or later}
\begin{document}
\maketitle
\setcounter{subsection}{1}
\noindent\textit{Editorial note.} All delivered mathematical text is the author's own, copied from the pinned Stacks Project chapter file cotangent.tex; no statement or proof step was rewritten, reordered or removed.\par\medskip
\noindent\textit{Editorial note.} The statement displays the triangle; the proof parallels the topoi case with the structure sheaves of the ringed spaces.\par\medskip
\noindent\textit{Editorial note.} The two references of the counted spans are recorded as non-load-bearing navigation pointers; the argument that consumes them stays inside the counted proof span.\par\medskip
\noindent\textit{Editorial note.} Preamble: the retained preamble of this chapter (sources/stacks-cotangent/retained-preamble.tex, lines 1-63), containing the Xy-pic setup, the comment environments, the theorem declarations and the macros including \textbackslash{}NL.\par\medskip
\noindent\textit{Allowed standard prerequisites (editorial).}\begin{itemize}
\item The pullback description of the cotangent complex. Lemma lemma-pullback-cotangent-morphism-topoi, used at proof lines 3773-3777. Identifies the pullback term of the triangle.
\item The distinguished triangle for a composition of sheaves of rings. Lemma lemma-triangle-sheaves-rings (3494-3505), used at proof lines 3778-3782. Is the sheaf-theoretic input instantiated by the three structure sheaves.
\item The definition of the cotangent complex of a morphism of ringed spaces. used at proof lines 3765-3770. Fixes the objects appearing in the displayed triangle.
\end{itemize}\medskip
\ifdefined\hypertarget\hypertarget{sourcepointer-1}{}\fi
\noindent\textit{Source reference.} Original lemma 08SV: focused-seeds/AGP/sources/stacks-snapshot-20261011/tags.txt, source lines 11143--11143. Navigation pointer into the Stacks Project; the standard fact it supplies is applied inside the counted proof span and the argument resting on it stays there.\par\medskip
\ifdefined\hypertarget\hypertarget{sourcepointer-2}{}\fi
\noindent\textit{Source reference.} Original lemma 08SZ: focused-seeds/AGP/sources/stacks-snapshot-20261011/tags.txt, source lines 11147--11147. Navigation pointer into the Stacks Project; the standard fact it supplies is applied inside the counted proof span and the argument resting on it stays there.\par\medskip
\begin{lemma}
\label{lemma-triangle-ringed-spaces}
Let $f : X \to Y$ and $g : Y \to Z$ be morphisms of ringed spaces.
Then there is a canonical distinguished triangle
$$
Lf^* L_{Y/Z} \to L_{X/Z} \to L_{X/Y} \to Lf^*L_{Y/Z}[1]
$$
in $D(\mathcal{O}_X)$.
\end{lemma}

\begin{proof}
Set $h = g \circ f$ so that $h^{-1}\mathcal{O}_Z = f^{-1}g^{-1}\mathcal{O}_Z$.
By Lemma \ref{lemma-pullback-cotangent-morphism-topoi} we have
$f^{-1}L_{Y/Z} = L_{f^{-1}\mathcal{O}_Y/h^{-1}\mathcal{O}_Z}$
and this is a complex of flat $f^{-1}\mathcal{O}_Y$-modules.
Hence the distinguished triangle above is an example of the
distinguished triangle of
Lemma \ref{lemma-triangle-sheaves-rings}
with $\mathcal{A} = h^{-1}\mathcal{O}_Z$, $\mathcal{B} = f^{-1}\mathcal{O}_Y$,
and $\mathcal{C} = \mathcal{O}_X$.
\end{proof}

\end{document}
